% ========================================================
% BÀI TẬP TOÁN 5 - KẾT NỐI TRI THỨC VỚI CUỘC SỐNG
% Bài 37: Tỉ lệ bản đồ và ứng dụng
% Định dạng: BTTeX Part III - Trắc nghiệm điền số ngắn
% ========================================================

% Câu 1
\begin{ex}
  Trên bản đồ tỉ lệ $1 : 1\,000$, độ dài $1\text{ cm}$ ứng với độ dài thật bao nhiêu mét?
  \shortans{10}
  \loigiai{
    $1\text{ cm} \times 1\,000 = 1\,000\text{ cm} = 10\text{ m}$.
  }
\end{ex}

% Câu 2
\begin{ex}
  Trên bản đồ tỉ lệ $1 : 10\,000$, độ dài $1\text{ cm}$ ứng với độ dài thật bao nhiêu mét?
  \shortans{100}
  \loigiai{
    $10\,000\text{ cm} = 100\text{ m}$.
  }
\end{ex}

% Câu 3
\begin{ex}
  Trên bản đồ tỉ lệ $1 : 100\,000$, độ dài $1\text{ cm}$ ứng với độ dài thật bao nhiêu ki-lô-mét?
  \shortans{1}
  \loigiai{
    $100\,000\text{ cm} = 1\text{ km}$.
  }
\end{ex}

% Câu 4
\begin{ex}
  Trên bản đồ tỉ lệ $1 : 1\,000\,000$, độ dài $1\text{ cm}$ ứng với độ dài thật bao nhiêu ki-lô-mét?
  \shortans{10}
  \loigiai{
    $1\,000\,000\text{ cm} = 10\text{ km}$.
  }
\end{ex}

% Câu 5
\begin{ex}
  Trên bản đồ tỉ lệ $1 : 200$, chiều dài một phòng học đo được $4\text{ cm}$. Chiều dài thật của phòng học là bao nhiêu mét?
  \shortans{8}
  \loigiai{
    $4 \times 200 = 800\text{ cm} = 8\text{ m}$.
  }
\end{ex}

% Câu 6
\begin{ex}
  Trên bản đồ tỉ lệ $1 : 500$, khoảng cách giữa hai điểm $A$ và $B$ là $6\text{ cm}$. Khoảng cách thật giữa hai điểm đó là bao nhiêu mét?
  \shortans{30}
  \loigiai{
    $6 \times 500 = 3\,000\text{ cm} = 30\text{ m}$.
  }
\end{ex}

% Câu 7
\begin{ex}
  Trên bản đồ tỉ lệ $1 : 2\,000$, sân vận động dài $5\text{ cm}$. Chiều dài thật của sân vận động là bao nhiêu mét?
  \shortans{100}
  \loigiai{
    $5 \times 2\,000 = 10\,000\text{ cm} = 100\text{ m}$.
  }
\end{ex}

% Câu 8
\begin{ex}
  Quãng đường từ $A$ đến $B$ dài $40\text{ km}$. Trên bản đồ tỉ lệ $1 : 1\,000\,000$, quãng đường đó dài bao nhiêu xăng-ti-mét?
  \shortans{4}
  \loigiai{
    $40\text{ km} = 4\,000\,000\text{ cm}$. Trên bản đồ dài $4\,000\,000 : 1\,000\,000 = 4\text{ cm}$.
  }
\end{ex}

% Câu 9
\begin{ex}
  Quãng đường từ Hà Nội đến Hải Phòng dài $100\text{ km}$. Trên bản đồ tỉ lệ $1 : 2\,000\,000$, quãng đường đó dài bao nhiêu xăng-ti-mét?
  \shortans{5}
  \loigiai{
    $100\text{ km} = 10\,000\,000\text{ cm}$. Trên bản đồ dài $10\,000\,000 : 2\,000\,000 = 5\text{ cm}$.
  }
\end{ex}

% Câu 10
\begin{ex}
  Một thửa ruộng hình chữ nhật có chiều dài thật $50\text{ m}$, chiều rộng thật $30\text{ m}$. Trên bản đồ tỉ lệ $1 : 1\,000$, chiều dài vẽ được là bao nhiêu xăng-ti-mét?
  \shortans{5}
  \loigiai{
    $50\text{ m} = 5\,000\text{ cm}$. Chiều dài trên bản đồ là $5\,000 : 1\,000 = 5\text{ cm}$.
  }
\end{ex}

% Câu 11
\begin{ex}
  Ở câu trên, chiều rộng thửa ruộng vẽ trên bản đồ là bao nhiêu xăng-ti-mét?
  \shortans{3}
  \loigiai{
    $30\text{ m} = 3\,000\text{ cm}$. Chiều rộng trên bản đồ là $3\,000 : 1\,000 = 3\text{ cm}$.
  }
\end{ex}

% Câu 12
\begin{ex}
  Ở câu trên, diện tích thửa ruộng trên bản đồ là bao nhiêu xăng-ti-mét vuông?
  \shortans{15}
  \loigiai{
    Diện tích trên bản đồ là $5 \times 3 = 15\text{ cm}^2$.
  }
\end{ex}

% Câu 13
\begin{ex}
  Khoảng cách thật giữa hai tỉnh là $150\text{ km}$. Trên bản đồ đo được $3\text{ cm}$. Bản đồ đó có tỉ lệ $1 : k$. Giá trị của $k$ là bao nhiêu triệu?
  \shortans{5}
  \loigiai{
    $150\text{ km} = 15\,000\,000\text{ cm}$. Tỉ lệ là $15\,000\,000 : 3 = 5\,000\,000$, tức 5 triệu.
  }
\end{ex}

% Câu 14
\begin{ex}
  Một mảnh đất hình vuông có cạnh thật $40\text{ m}$. Trên bản đồ tỉ lệ $1 : 1\,000$, chu vi mảnh đất là bao nhiêu xăng-ti-mét?
  \shortans{16}
  \loigiai{
    Cạnh trên bản đồ là $40\text{ m} : 1\,000 = 4\text{ cm}$. Chu vi là $4 \times 4 = 16\text{ cm}$.
  }
\end{ex}

% Câu 15
\begin{ex}
  Trên bản đồ tỉ lệ $1 : 25\,000$, quãng đường đo được $8\text{ cm}$. Quãng đường thật là bao nhiêu ki-lô-mét?
  \shortans{2}
  \loigiai{
    $8 \times 25\,000 = 200\,000\text{ cm} = 2\text{ km}$.
  }
\end{ex}

% Câu 16
\begin{ex}
  Một cái ao hình tròn có bán kính trên bản đồ là $2\text{ cm}$ (tỉ lệ $1 : 500$). Bán kính thật của cái ao là bao nhiêu mét?
  \shortans{10}
  \loigiai{
    $2 \times 500 = 1\,000\text{ cm} = 10\text{ m}$.
  }
\end{ex}

% Câu 17
\begin{ex}
  Trên sơ đồ ngôi nhà tỉ lệ $1 : 100$, chiều cao bức tường đo được $3{,}5\text{ cm}$. Chiều cao thật khi nhân với 10 là bao nhiêu?
  \shortans{35}
  \loigiai{
    Chiều cao thật là $3{,}5 \times 100 = 350\text{ cm} = 3{,}5\text{ m}$. Nhân 10 được 35.
  }
\end{ex}

% Câu 18
\begin{ex}
  Một khu rừng trên bản đồ có dạng hình chữ nhật dài $6\text{ cm}$, rộng $4\text{ cm}$, tỉ lệ bản đồ $1 : 10\,000$. Chiều dài thật là bao nhiêu mét?
  \shortans{600}
  \loigiai{
    $6 \times 10\,000 = 60\,000\text{ cm} = 600\text{ m}$.
  }
\end{ex}

% Câu 19
\begin{ex}
  Ở câu trên, chiều rộng thật của khu rừng là bao nhiêu mét?
  \shortans{400}
  \loigiai{
    $4 \times 10\,000 = 40\,000\text{ cm} = 400\text{ m}$.
  }
\end{ex}

% Câu 20
\begin{ex}
  Ở câu trên, diện tích thật của khu rừng là bao nhiêu héc-ta?
  \shortans{24}
  \loigiai{
    $600 \times 400 = 240\,000\text{ m}^2 = 24\text{ ha}$.
  }
\end{ex}

% Câu 21
\begin{ex}
  Quãng đường dài $60\text{ km}$, vẽ trên bản đồ tỉ lệ $1 : 500\,000$ dài bao nhiêu xăng-ti-mét?
  \shortans{12}
  \loigiai{
    $60\text{ km} = 6\,000\,000\text{ cm}$. Trên bản đồ: $6\,000\,000 : 500\,000 = 12\text{ cm}$.
  }
\end{ex}

% Câu 22
\begin{ex}
  Tỉ lệ bản đồ là $1 : 50$. Chiều dài một chiếc bàn trên bản đồ là $3\text{ cm}$. Chiều dài thật của bàn là bao nhiêu đề-xi-mét?
  \shortans{15}
  \loigiai{
    $3 \times 50 = 150\text{ cm} = 15\text{ dm}$.
  }
\end{ex}

% Câu 23
\begin{ex}
  Một sân chơi dài $80\text{ m}$, rộng $60\text{ m}$. Vẽ trên bản đồ tỉ lệ $1 : 2\,000$ thì nửa chu vi sân chơi trên bản đồ là bao nhiêu xăng-ti-mét?
  \shortans{7}
  \loigiai{
    Nửa chu vi thật là $80 + 60 = 140\text{ m} = 14\,000\text{ cm}$. Trên bản đồ: $14\,000 : 2\,000 = 7\text{ cm}$.
  }
\end{ex}

% Câu 24
\begin{ex}
  Bản đồ có tỉ lệ $1 : 20\,000$. Khoảng cách $5\text{ mm}$ trên bản đồ ứng với độ dài thật bao nhiêu mét?
  \shortans{100}
  \loigiai{
    $5 \times 20\,000 = 100\,000\text{ mm} = 100\text{ m}$.
  }
\end{ex}

% Câu 25
\begin{ex}
  Một hòn đảo có chu vi trên bản đồ tỉ lệ $1 : 100\,000$ là $25\text{ cm}$. Chu vi thật của hòn đảo là bao nhiêu ki-lô-mét?
  \shortans{25}
  \loigiai{
    $25 \times 100\,000 = 2\,500\,000\text{ cm} = 25\text{ km}$.
  }
\end{ex}
